\documentclass[12pt,letterpaper]{article}

\usepackage[utf8]{inputenc}
\usepackage[T1]{fontenc}
\usepackage[spanish,es-noshorthands]{babel}
\usepackage{amsmath,amssymb,amsthm}
\usepackage{geometry}
\usepackage{enumitem}
\usepackage{xcolor}
\usepackage{hyperref}

\geometry{letterpaper, margin=1in}

\definecolor{JustifColor}{HTML}{3B4A66}
\definecolor{ResultColor}{HTML}{0E7490}

\newcommand{\justif}[1]{%
  \par\noindent
  {\small\color{JustifColor}%
  \begin{quote}
  \textit{Justificación.} #1
  \end{quote}}%
}

\newenvironment{resultado}
  {\begin{quote}\color{ResultColor}\bfseries}
  {\end{quote}}

\newcommand{\paso}[1]{\subsubsection*{#1}}

\hypersetup{
  colorlinks=true,
  linkcolor=ResultColor,
  citecolor=ResultColor,
  urlcolor=ResultColor,
  pdftitle={El Factor Integrante: Demostración Paso a Paso},
}

\title{\textbf{El Factor Integrante}\\[4pt]
\large Demostración paso a paso, con justificaciones, para ecuaciones lineales\\ y para forzar la exactitud de una ecuación diferencial}
\author{Cálculo III}
\date{}

\begin{document}
\maketitle
\thispagestyle{empty}

\noindent
Este documento reúne, con todo detalle, la construcción del factor integrante
en sus dos apariciones del curso: primero como la herramienta que resuelve
\emph{toda} ecuación lineal de primer orden (Sesión 4), y después como el
recurso que rescata una ecuación \emph{no exacta} (Sesión 5). Al final se
muestra que la primera es, de hecho, un caso particular de la segunda.

\tableofcontents

\bigskip
\hrule
\bigskip

\section{Factor integrante para ecuaciones lineales}

\subsection*{Planteamiento}

Toda ecuación diferencial lineal de primer orden se puede escribir en su
\textbf{forma estándar}:
\begin{equation}
\frac{dy}{dx} + P(x)\,y = f(x), \tag{1}
\end{equation}
con $P$ y $f$ continuas en un intervalo $I$. El lado izquierdo de (1) no es,
en general, la derivada de ningún producto simple. La idea del factor
integrante es \emph{multiplicar toda la ecuación por una función bien
elegida} para que, después de multiplicar, sí lo sea.

\paso{Paso 1 --- El ansatz}

Buscamos una función $\mu(x)$, distinta de cero en $I$, tal que
\begin{equation}
\mu(x)\left[\frac{dy}{dx} + P(x)\,y\right] \;=\; \frac{d}{dx}\Big[\mu(x)\,y\Big]. \tag{2}
\end{equation}

\justif{Si logramos construir tal $\mu$, el problema se reduce a una
integración directa: el lado derecho de (2) se integra sin esfuerzo respecto
de $x$, y de ahí despejamos $y$. Todo el método depende de que este $\mu$
exista y se pueda calcular explícitamente.}

\paso{Paso 2 --- Expandir el lado derecho}

Por la regla del producto,
\begin{equation}
\frac{d}{dx}\big[\mu(x)\,y\big] = \mu(x)\,\frac{dy}{dx} + \mu'(x)\,y. \tag{3}
\end{equation}

\paso{Paso 3 --- Igualar y comparar coeficientes}

Sustituyendo (3) en (2):
\[
\mu\,\frac{dy}{dx} + \mu P\, y \;=\; \mu\,\frac{dy}{dx} + \mu'\, y.
\]
El término $\mu\, dy/dx$ aparece idéntico en ambos lados y se cancela, dejando
\begin{equation}
\mu(x)\,P(x)\,y \;=\; \mu'(x)\,y. \tag{4}
\end{equation}

\justif{La igualdad (4) debe cumplirse para la solución $y(x)$ que
\emph{todavía no conocemos}: $\mu$ se construye \emph{antes} de resolver la
ecuación, sin usar ninguna información sobre $y$. Para que la identidad sea
válida sin importar cuál resulte ser $y$, los coeficientes que multiplican a
$y$ en ambos lados deben coincidir exactamente. Esto elimina a $y$ del
problema por completo y deja una ecuación solo para $\mu$:
\[
\mu'(x) = P(x)\,\mu(x).
\]}

\paso{Paso 4 --- Resolver la ecuación para $\mu$}

La ecuación $\mu' = P\mu$ es, ella misma, \textbf{separable}:
\[
\frac{d\mu}{\mu} = P(x)\,dx.
\]
Integrando ambos lados,
\[
\ln|\mu| = \int P(x)\,dx + C_1.
\]

\justif{Aquí se usa que $\mu \neq 0$ en $I$ (lo exigimos desde el Paso 1),
así que $1/\mu$ está bien definida y la separación de variables es válida.}

\paso{Paso 5 --- Despejar $\mu$ y elegir la constante}

Exponenciando,
\[
\mu(x) = e^{C_1}\, e^{\int P(x)\,dx}.
\]

\justif{Cualquier múltiplo constante de un factor integrante es también un
factor integrante: si $\mu$ satisface (2), entonces $k\mu$ lo satisface
igual para cualquier constante $k \neq 0$, porque $k$ se cancela en la
ecuación (4) del Paso 3. Por lo tanto la constante $C_1$ es enteramente una
cuestión de conveniencia, no de necesidad matemática. Elegimos la opción más
simple, $C_1 = 0$ (equivalentemente $e^{C_1}=1$):}

\begin{resultado}
\[
\mu(x) = e^{\int P(x)\,dx}
\]
\end{resultado}

\subsection*{Comprobación de que el método funciona}

Con este $\mu$, la ecuación (2) se cumple por construcción, así que
multiplicar (1) por $\mu(x)$ da directamente
\[
\frac{d}{dx}\Big[\mu(x)\,y\Big] = \mu(x)\,f(x).
\]
Integrando respecto de $x$:
\[
\mu(x)\,y = \int \mu(x)\,f(x)\,dx + C,
\]
\begin{equation}
y = \frac{1}{\mu(x)}\left[\int \mu(x)\,f(x)\,dx + C\right]. \tag{5}
\end{equation}

\justif{La división por $\mu(x)$ en (5) siempre es válida: como
$\mu(x)=e^{\int P\,dx}$ es una función exponencial, se tiene $\mu(x)>0$ para
\emph{todo} $x \in I$ --- nunca se anula. Este es precisamente el papel que
le pedimos a $\mu$ desde el Paso 1, y ahora queda garantizado
automáticamente por la fórmula que obtuvimos.}

\bigskip
\hrule
\bigskip

\section{Factor integrante para volver exacta una ecuación}

\subsection*{Planteamiento}

Sea la ecuación diferencial en forma diferencial
\[
M(x,y)\,dx + N(x,y)\,dy = 0,
\]
que \textbf{no} es exacta, es decir, $\dfrac{\partial M}{\partial y} \neq
\dfrac{\partial N}{\partial x}$. Buscamos una función $\mu$ tal que
\begin{equation}
\mu(x,y)\,M(x,y)\,dx + \mu(x,y)\,N(x,y)\,dy = 0 \tag{6}
\end{equation}
sí sea exacta.

\paso{Paso 1 --- La condición general de exactitud para la nueva ecuación}

Por el criterio de exactitud aplicado a (6), necesitamos
\[
\frac{\partial}{\partial y}\big[\mu M\big] = \frac{\partial}{\partial x}\big[\mu N\big].
\]
Expandiendo ambos lados con la regla del producto --- ahora $\mu = \mu(x,y)$
depende, en principio, de ambas variables ---:
\[
\mu_y M + \mu M_y = \mu_x N + \mu N_x.
\]
Reacomodando:
\begin{equation}
M\,\mu_y - N\,\mu_x = \mu\,(N_x - M_y). \tag{$\ast$}
\end{equation}

\justif{La ecuación $(\ast)$ es una \textbf{ecuación diferencial parcial}
para $\mu(x,y)$, de primer orden y, en general, tan difícil de resolver como
la ecuación original (o más). Por eso no la atacamos en su forma más
general: buscamos soluciones particulares bajo un supuesto simplificador
--- un \emph{ansatz} --- que la reduzca a una ecuación diferencial
\textbf{ordinaria}, mucho más manejable. Los dos ansatz más útiles son
suponer que $\mu$ depende de una sola variable.}

\paso{Caso A --- $\mu = \mu(x)$ (depende solo de $x$)}

Si $\mu$ no depende de $y$, entonces $\mu_y = 0$, y $(\ast)$ se simplifica a
\[
-N\,\mu_x = \mu\,(N_x - M_y)
\quad\Longrightarrow\quad
\frac{\mu_x}{\mu} = \frac{M_y - N_x}{N}.
\]

\justif{El lado izquierdo, $\mu_x/\mu = \mu'(x)/\mu(x)$, depende solo de
$x$ por construcción --- es la derivada logarítmica de una función de una
sola variable. Para que la igualdad tenga sentido, el lado derecho,
$(M_y-N_x)/N$, debe depender \emph{también} solo de $x$, aunque $M$, $N$,
$M_y$ y $N_x$ por separado dependan de $x$ y de $y$. \textbf{Esta es
precisamente la condición práctica que se verifica antes de usar la
fórmula}: si al calcular $(M_y-N_x)/N$ toda la $y$ se cancela, el Caso A
aplica.}

Cuando esa condición se cumple, queda una ecuación separable ordinaria para
$\mu$ --- idéntica en forma a la del Paso 4 de la Sección 1 ---:
\[
\frac{d\mu}{\mu} = \frac{M_y-N_x}{N}\,dx.
\]

\begin{resultado}
\[
\mu(x) = e^{\int \frac{M_y-N_x}{N}\,dx}
\]
\end{resultado}

\paso{Caso B --- $\mu = \mu(y)$ (depende solo de $y$)}

Simétricamente, si $\mu_x = 0$, la ecuación $(\ast)$ se simplifica a
\[
M\,\mu_y = \mu\,(N_x - M_y)
\quad\Longrightarrow\quad
\frac{\mu_y}{\mu} = \frac{N_x - M_y}{M}.
\]

\justif{Mismo argumento que en el Caso A, intercambiando los papeles de $x$
y $y$: la igualdad solo es consistente si $(N_x-M_y)/M$ depende únicamente
de $y$.}

\begin{resultado}
\[
\mu(y) = e^{\int \frac{N_x-M_y}{M}\,dy}
\]
\end{resultado}

\subsection*{Comprobación de que el método funciona}

\textbf{Caso A.} Sea $\mu=\mu(x)$ construido como arriba, y definamos
$M^{*} = \mu M$, $N^{*} = \mu N$. Queremos ver que $M^{*}_y = N^{*}_x$:
\[
M^{*}_y = \mu\, M_y,
\qquad
N^{*}_x = \mu'\,N + \mu\, N_x.
\]
Usando $\mu' = \mu \cdot \dfrac{M_y-N_x}{N}$ en la segunda expresión:
\[
N^{*}_x
= \mu\cdot\frac{M_y-N_x}{N}\cdot N + \mu\,N_x
= \mu\,(M_y-N_x) + \mu\,N_x
= \mu\,M_y
= M^{*}_y. \checkmark
\]

\textbf{Caso B.} De manera simétrica, con $M^{*}_y = \mu' M + \mu M_y$ y
$N^{*}_x = \mu N_x$, y usando $\mu' = \mu\cdot\dfrac{N_x-M_y}{M}$:
\[
M^{*}_y
= \mu\cdot\frac{N_x-M_y}{M}\cdot M + \mu\,M_y
= \mu\,(N_x-M_y) + \mu\,M_y
= \mu\,N_x
= N^{*}_x. \checkmark
\]

\justif{En ambos casos, multiplicar por el $\mu$ construido garantiza
exactitud \emph{por construcción}, no por coincidencia: la cancelación
algebraica funciona exactamente porque $\mu$ fue diseñado en el Paso 1 de
cada caso para lograrlo.}

\bigskip
\hrule
\bigskip

\section{Observación final: el factor lineal es un caso particular}

Las Secciones 1 y 2 no son dos técnicas independientes. La ecuación lineal
\[
\frac{dy}{dx} + P(x)\,y = f(x)
\]
se puede reescribir en forma diferencial como
\[
\big[P(x)\,y - f(x)\big]\,dx + dy = 0,
\]
es decir, con $M(x,y) = P(x)\,y - f(x)$ y $N(x,y) = 1$. Entonces
\[
M_y = P(x), \qquad N_x = 0,
\]
\[
\frac{M_y - N_x}{N} = \frac{P(x) - 0}{1} = P(x),
\]
que depende solo de $x$ \textbf{automáticamente}, para cualquier ecuación
lineal, sin excepción. Aplicando la fórmula del Caso A de la Sección 2:
\[
\mu(x) = e^{\int P(x)\,dx},
\]
exactamente el factor integrante de la Sección 1.

\begin{resultado}
El método de la Sección 1 no es una técnica aparte: es el Caso A del factor
integrante para ecuaciones exactas, aplicado a una ecuación que --- por su
propia estructura lineal --- siempre cae en ese caso.
\end{resultado}

\end{document}
